\documentclass[12pt,a4paper]{article}
\usepackage[margin=25mm, headheight=15pt, headsep=10pt]{geometry}
\usepackage{fontspec}
\usepackage[table]{xcolor} % Note: table option is required for rowcolors
\usepackage{titlesec}
\usepackage{tocloft}
\usepackage{fancyhdr}
\usepackage{booktabs}
\usepackage{tcolorbox}
\tcbuselibrary{skins, breakable}
\usepackage[colorlinks=true, linkcolor=emerald, urlcolor=emerald, bookmarks=true]{hyperref}

\definecolor{emerald}{RGB}{0,102,51}
\definecolor{darkred}{RGB}{139,0,0}
\definecolor{lightgray}{RGB}{245,245,245}

\usepackage[nil]{babel}
\babelprovide[import, main]{bengali}
\babelprovide[import]{arabic}
\babelfont{rm}[Renderer=HarfBuzz,Scale=1.0]{Noto Serif Bengali}
\babelfont{sf}[Renderer=HarfBuzz]{Noto Sans Bengali}
\babelfont{tt}[Renderer=HarfBuzz]{Noto Sans Bengali}
\babelfont[arabic]{rm}[Renderer=HarfBuzz,Scale=1.1]{Noto Naskh Arabic}

% Section styling
\titleformat{\section}{\Large\bfseries\color{emerald}}{\thesection.}{0.5em}{}
\titleformat{\subsection}{\large\bfseries\color{emerald!85!black}}{\thesubsection.}{0.5em}{}
\titleformat{\subsubsection}{\normalsize\bfseries\color{emerald!70!black}}{\thesubsubsection.}{0.5em}{}
\titlespacing*{\section}{0pt}{18pt}{8pt}

% Fancy Header/Footer setup
\pagestyle{fancy}
\fancyhf{} % clear all header and footer fields
\fancyhead[L]{\color{emerald}\small\textbf{\nouppercase{\leftmark}}}
\fancyfoot[L]{\scriptsize\color{black!50}{\lat{AI}}-সংকলিত, আলেম-যাচাইকৃত নয়}
\fancyfoot[C]{\color{emerald!70!black}\thepage}
\renewcommand{\headrulewidth}{0.4pt}
\renewcommand{\headrule}{\hbox to\headwidth{\color{emerald}\leaders\hrule height \headrulewidth\hfill}}
\renewcommand{\footrulewidth}{0pt}

% Custom boxes
% 1. Ayat/Hadith Box
\newtcolorbox{ayatbox}[1][]{
    enhanced, breakable, 
    colback=emerald!5, 
    colframe=emerald, 
    boxrule=0.5pt, 
    arc=3pt, 
    left=10pt, right=10pt, top=10pt, bottom=10pt,
    #1
}

% 2. Quote/Translation Box
\newtcolorbox{quotebox}[1][]{
    enhanced, breakable,
    frame hidden,
    colback=lightgray,
    borderline west={3pt}{0pt}{emerald},
    left=15pt, right=10pt, top=10pt, bottom=10pt,
    sharp corners,
    #1
}

% 3. AI-generated notice box
\newtcolorbox{warnbox}[1][]{
    enhanced, breakable,
    colback=red!4,
    colframe=darkred,
    boxrule=0.8pt,
    arc=3pt,
    left=10pt, right=10pt, top=10pt, bottom=10pt,
    #1
}

% Commands
\newcommand{\arab}[1]{\foreignlanguage{arabic}{#1}}
\newcommand{\kib}[1]{\textcolor{emerald}{\textbf{#1}}}

% Latin (English) fallback: Noto Serif Bengali has NO Latin glyphs
\newfontfamily\latfont[Renderer=HarfBuzz]{Noto Serif}
\newcommand{\lat}[1]{{\latfont #1}}

\setlength{\parskip}{5pt}
\setlength{\parindent}{0pt}

% Fix for missing bullet in Noto Serif Bengali
\renewcommand{\labelitemi}{$\bullet$}



\begin{document}
\begin{center}{\Large\bfseries\color{emerald} \lat{10}-তাশাহহুদ-প্রথম}\end{center}
\vspace{1em}
\section*{১০ — তাশাহহুদ (প্রথম বৈঠক): \textit{\lat{Attahiyyatu}...}}

\begin{warnbox}
 \textbf{গুরুত্বপূর্ণ নোটিশ:} এই কন্টেন্টটি \lat{AI} দিয়ে প্রস্তুত। কেবল সহীহ/হাসান হাদিস রাখা হয়েছে। আলেম-যাচাইকৃত নয়।
\end{warnbox}

\medskip\hrule\medskip

\subsection*{১. সূত্র কার্ড}

\begin{table}[h]\centering\small
\begin{tabular}{ll}
বিষয় & বিবরণ \\
\hline
\textbf{বিষয়} & প্রথম তাশাহহুদ — ২/৩/৪ রাকাতের পর বসে পড়া হয় \\
\textbf{সূত্র} & সহীহ বুখারি ৬২৬৫, সহীহ মুসলিম ৪০২ \\
\textbf{মর্যাদা} & \textbf{ফরজ} — এটি নামাজের অবিচ্ছেদ্য অংশ \\
\hline
\end{tabular}
\end{table}

\medskip\hrule\medskip

\subsection*{২. মূল আরবি}

\begin{center}\arab{التَّحِيَّاتُ لِلَّهِ وَالصَّلَوَاتُ وَالطَّيِّبَاتُ، السَّلَامُ عَلَيْكَ أَيُّهَا النَّبِيُّ وَرَحْمَةُ اللَّهِ وَبَرَكَاتُهُ، السَّلَامُ عَلَيْنَا وَعَلَىٰ عِبَادِ اللَّهِ الصَّالِحِينَ، أَشْهَدُ أَنْ لَا إِلَٰهَ إِلَّا اللَّهُ وَأَشْهَدُ أَنَّ مُحَمَّدًا عَبْدُهُ وَرَسُولُهُ}\end{center}

\medskip\hrule\medskip

\subsection*{৩. বাংলা উচ্চারণ}

\textbf{আততাহিয়্যাতু লিল্লাহি ওয়াস সালাওয়াতু ওয়াঃ তাইয়্যিবাত, আসসালামু আলাইকা আইয়ুহান নাবিয়্যু ওয়া রাহমাতুল্লাহি ওয়া বারাকাতুহ, আসসালামু আলাইনা ওয়া আলা ইবাদিল্লাহিস সালিহীন, আশহাদু আল্লা ইলাহা ইল্লাল্লাহু ওয়া আশহাদু আন্না মুহাম্মাদান আবদুহু ওয়া রাসূলুহ}

\medskip\hrule\medskip

\subsection*{৪. শব্দ-শব্দ অর্থ}

\begin{table}[h]\centering\small
\begin{tabular}{lll}
আরবি & বাংলা & \lat{English} \\
\hline
\textbf{\arab{التَّحِيَّاتُ}} & সালাম/প্রণাম & \textit{\lat{Greetings}} \\
\textbf{\arab{لِلَّهِ}} & আল্লাহর জন্য & \textit{\lat{are} \lat{for} \lat{Allah}} \\
\textbf{\arab{وَالصَّلَوَاتُ}} & ও দরুদ/নামাজ & \textit{\lat{and} \lat{prayers}} \\
\textbf{\arab{وَالطَّيِّبَاتُ}} & ও পবিত্র বাক্য & \textit{\lat{and} \lat{pure} \lat{words}} \\
\textbf{\arab{السَّلَامُ} \arab{عَلَيْكَ}} & তোমার ওপর শান্তি & \textit{\lat{Peace} \lat{be} \lat{upon} \lat{you}} \\
\textbf{\arab{أَيُّهَا} \arab{النَّبِيُّ}} & হে নবী & \textit{\lat{O} \lat{Prophet}} \\
\textbf{\arab{وَرَحْمَةُ} \arab{اللَّهِ}} & ও আল্লাহর দয়া & \textit{\lat{and} \lat{Allah's} \lat{mercy}} \\
\textbf{\arab{وَبَرَكَاتُهُ}} & ও তাঁর বরকত & \textit{\lat{and} \lat{His} \lat{blessings}} \\
\textbf{\arab{السَّلَامُ} \arab{عَلَيْنَا}} & আমাদের ওপর শান্তি & \textit{\lat{Peace} \lat{be} \lat{upon} \lat{us}} \\
\textbf{\arab{وَعَلَىٰ} \arab{عِبَادِ} \arab{اللَّهِ}} & ও আল্লাহর বান্দাদের ওপর & \textit{\lat{and} \lat{upon} \lat{Allah's} \lat{servants}} \\
\textbf{\arab{الصَّالِحِينَ}} & সৎকর্মশীলদের & \textit{\lat{the} \lat{righteous}} \\
\textbf{\arab{أَشْهَدُ} \arab{أَنْ} \arab{لَا} \arab{إِلَٰهَ}} & আমি সাক্ষ্য দিই কোনো ইলাহ নেই & \textit{\lat{I} \lat{bear} \lat{witness} \lat{there} \lat{is} \lat{no} \lat{god}} \\
\textbf{\arab{إِلَّا} \arab{اللَّهُ}} & আল্লাহ ছাড়া & \textit{\lat{except} \lat{Allah}} \\
\textbf{\arab{وَأَشْهَدُ} \arab{أَنَّ}} & ও আমি সাক্ষ্য দিই & \textit{\lat{and} \lat{I} \lat{bear} \lat{witness}} \\
\textbf{\arab{مُحَمَّدًا}} & মুহাম্মদ & \textit{\lat{that} \lat{Muhammad}} \\
\textbf{\arab{عَبْدُهُ}} & তাঁর বান্দা & \textit{\lat{is} \lat{His} \lat{servant}} \\
\textbf{\arab{وَرَسُولُهُ}} & ও তাঁর রাসূল & \textit{\lat{and} \lat{His} \lat{Messenger}} \\
\hline
\end{tabular}
\end{table}

\medskip\hrule\medskip

\subsection*{৫. পূর্ণ বাংলা অনুবাদ}

\textbf{\lat{“}সকল প্রণাম আল্লাহর জন্য, সকল দরুদ ও পবিত্র বাক্য তাঁর জন্য। হে নবী! তোমার ওপর আল্লাহর শান্তি, দয়া ও বরকত বর্ষিত হোক। আমাদের ওপর ও আল্লাহর সৎকর্মশীল বান্দাদের ওপর শান্তি বর্ষিত হোক। আমি সাক্ষ্য দিই: আল্লাহ ছাড়া কোনো ইলাহ নেই এবং আমি সাক্ষ্য দিই মুহাম্মদ তাঁর বান্দা ও রাসূল।\lat{”}}

\medskip\hrule\medskip

\subsection*{৬. সংক্ষিপ্ত ব্যাখ্যা}

\subsubsection*{\textbf{"তাহিয়্যাত" — আল্লাহর প্রণাম}}
"\lat{Attahiyyat}" মানে সকল প্রণাম — জীবিত ও মৃত সকল সৃষ্টির প্রণাম আল্লাহর জন্য। এটি \textbf{\lat{tashahhud} (সাক্ষ্যদান)} — নামাজে বসে থাকার সময় আল্লাহর সামনে নিজেকে \textbf{\lat{humble} (বিনীত)} করা।

\subsubsection*{\textbf{"আশহাদু আল্লা ইলাহা ইল্লাল্লাহ" — কালেমা}}
এই অংশে \textbf{\lat{index} (আঙুল)} উঁচু করা হয় (শাফি'ি মাযহাব) — কারণ এটি \textbf{\lat{shahadah} (সাক্ষ্যদান)}। আঙুল উঁচু করা হলো \textbf{\lat{imony} (সাক্ষ্য)}-র প্রতীক।

\subsubsection*{\textbf{কীভাবে বসবেন?}}
\begin{itemize}
\item \textbf{বাম পা} বসার ভিত্তি — \textbf{ডান পা} উঁচু করে (পায়ের ডান পাশে বসুন)।
\item \textbf{হাত হাঁটুর ওপর} — দুই হাতের হাতল।
\item \textbf{নজর} — চোখ হাতের দিকে বা আঙুলে।
\end{itemize}
\medskip\hrule\medskip

\subsection*{৭. খুশুর জন্য মূল পয়েন্টস}

\begin{table}[h]\centering\small
\begin{tabular}{ll}
পয়েন্ট & কীভাবে প্রয়োগ করবেন \\
\hline
\textbf{গভীরতা} & প্রতিটি শব্দের অর্থ মনে করে পড়ুন \\
\textbf{দরুদ} & "আসসালামু আলাইকা আইয়ুহান নাবিয়্যু" — নবীজির প্রতি সালাম \\
\textbf{সাক্ষ্য} & "আশহাদু আল্লা ইলাহা ইল্লাল্লাহ" — দৃঢ় স্বরে বলুন \\
\textbf{ইতমিনান} & তাড়াহুড়ো নয় — পূর্ণ শান্তিতে বসুন \\
\hline
\end{tabular}
\end{table}

\medskip\hrule\medskip

\subsection*{৮. সংশ্লিষ্ট হাদিস}

\begin{table}[h]\centering\small
\begin{tabular}{ll}
হাদিস & গ্রন্থ \\
\hline
\textbf{\lat{“}তোমাদের প্রত্যেকে তাশাহহুদ পড়ুক।\lat{”}} & সহীহ বুখারি ৬২৬৫ \\
\textbf{\lat{“}নামাজে তাশাহহুদ ফরজ।\lat{”}} & সহীহ মুসলিম ৪০২ \\
\hline
\end{tabular}
\end{table}

\medskip\hrule\medskip

\textit{শেষ আপডেট: আগস্ট ২০২৬}

\end{document}
